\section{Groups, modules, and dimension as instances}\label{e:section}

The comparison theorem has been proved before choosing a multiplication
table, a scalar field, or a basis. To use it, we must now check its hypotheses.
This is an essential part of the argument: a category name does not by itself
verify an image factorization or construct a quotient.

\subsection{One verification for algebraic structures}\label{e:varieties}

An algebraic theory specifies operations of finite arity and equations between
expressions built from them. A model is a set with those operations satisfying
the equations; a morphism preserves every operation. A subalgebra is a subset
closed under them. An internal relation on a model $A$ is a subalgebra of
$A\times A$. An equivalence relation of this kind is a \emph{congruence}:
equivalent inputs give equivalent outputs for every operation.

\begin{proposition}\label{e:algebraic-instances}
The category of models of an algebraic theory is Barr-exact. If its one-element
model is also initial, it is pointed. If the theory has a ternary term $p$
satisfying
\[
p(x,y,y)=x,\qquad p(x,x,z)=z,
\]
then it is a Mal'cev category. Under these last two conditions the general
comparison theorem applies to its finite composition series.
\end{proposition}
\begin{proof}
Products carry the operations coordinatewise. An equalizer is the subset
where two homomorphisms agree, closed under every operation. These give finite
limits, with the one-element algebra terminal. The image of a homomorphism
is a subalgebra, so every map factors as a surjection followed by an inclusion.
A pullback of a surjection is surjective: an element of the other object has
a lift of its image, giving a pair in the pullback.

Let $q:A\to B$ be surjective. A homomorphism $h:A\to C$ constant on its
fibers defines $\bar h(b)=h(a)$ for any $a$ with $q(a)=b$. This is independent
of $a$, and choosing lifts of the finitely many inputs of an operation proves
that $\bar h$ preserves that operation. Surjectivity proves uniqueness.
Thus $q$ is the coequalizer of its kernel pair.

Conversely a regular epimorphism is surjective. Factor a coequalizer
$q:A\to B$ through its set-theoretic image as $q=me$. If $q$ coequalizes
$u,v$, injectivity of $m$ shows $eu=ev$. Its universal property supplies
$d:B\to\im(q)$ with $dq=e$. Then $mdq=q$, so $md=\id_B$ because
coequalizers are epimorphisms. Also $dme=e$, and the surjectivity of $e$
gives $dm=\id_{\im(q)}$. Hence $m$ is an isomorphism.
We have established regularity with exactly the stated quotient maps.
Monomorphisms are precisely the injective homomorphisms: injectivity gives
left cancellation, while factoring a monomorphism into its image makes
the surjective part both monic and regular epic, hence an isomorphism by
Lemma~\ref{f:descent}. Thus the internal subobjects and relations really
are the subalgebras used in the next step.

For a congruence $E$ on $A$, form the set $A/E$ and define every operation
on classes by applying it to representatives. Closure of $E$ under the
operation makes this well-defined; the equations descend because they hold
on representatives. The quotient map has kernel pair exactly $E$ and has
the same coequalizer property just proved. Thus every internal equivalence
relation is effective, proving Barr-exactness.

Finally let $R\subseteq A\times A$ be reflexive and an internal relation.
It is closed under the term $p$. If $xRy$, apply $p$ to
$(y,y),(x,y),(x,x)$ to obtain $(y,x)$, proving symmetry. If $xRy$ and $yRz$,
apply $p$ to $(x,y),(y,y),(y,z)$ to obtain $(x,z)$, proving transitivity.
Thus every reflexive internal relation is an equivalence relation.
The pointed assertion is precisely the extra initial-object hypothesis.
\end{proof}

The displayed terms are operations in a presentation. For a nonabelian group,
$G^3\to G$, $(x,y,z)\mapsto xy^{-1}z$, need not be a group homomorphism.
What we used is that a subalgebra is closed under terms, including that term.
This distinction lets algebraic expressions justify a categorical property
without confusing them with arrows in the category.

\subsection{Reading the axioms in familiar examples}\label{e:examples}

For groups use multiplication, inverse and identity, with their axioms.
The one-element group is initial and terminal, and
$p(x,y,z)=xy^{-1}z$ has the required identities. For left modules over a fixed
unital ring $R$, use addition, negative, zero, and an operation $x\mapsto rx$
for every $r\in R$, with the module equations. The zero module is initial
and terminal, and $p(x,y,z)=x-y+z$ works. Neither argument uses a dimension,
a composition series, or Jordan--H\"older.

We check that the categorical quotient data are the familiar ones. In a
group $G$, the kernel $N$ of a homomorphism is normal, since
$q(gng^{-1})=q(g)q(n)q(g)^{-1}=1$. Conversely, for a normal subgroup $N$,
the relation $x\sim y$ if $x^{-1}y\in N$ is an equivalence relation
compatible with multiplication and inverses. For example, if $x'=xn$ and
$y'=ym$, then $x'y'=xy(y^{-1}ny)m$ differs from $xy$ by an element of $N$.
The congruence quotient just constructed is the usual group $G/N$.

More generally, if $E$ is a group congruence, its identity class $N$ is a
normal subgroup: closure gives products and inverses, and conjugating a pair
$(n,1)\in E$ by $(g,g)\in E$ gives $(gng^{-1},1)\in E$. Translating
$(x,y)$ by $(x^{-1},x^{-1})$ shows $xEy$ exactly when $x^{-1}y\in N$.
Thus normal subgroups and these quotient relations agree.

For modules, a congruence has zero class a submodule $N$, and translating
pairs gives $x\sim y$ exactly when $y-x\in N$. The quotient operations are
$(x+N)+(y+N)=x+y+N$ and $r(x+N)=rx+N$. Consequently all isomorphisms supplied
by the categorical theorem are $R$-linear. A quotient-simple group is exactly
a nontrivial group with no proper nontrivial normal subgroup; a quotient-simple
module is exactly a nonzero module with no proper nonzero submodule.

We therefore obtain both classical theorems from Theorem~\ref{j:jordan-holder}:
two composition series of a group have the same group factors, and two of an
$R$-module have the same module factors. This is one specialization step for
each category, not a second proof of the comparison theorem.

The same verification covers groups with specified endomorphism operators:
include those operators and their preservation equations in the theory.
Subobjects and quotients retain them. It also covers Lie algebras over a field,
using the additive term $x-y+z$, and nonunital associative algebras over a
commutative ring. Their kernels are respectively ideals and two-sided ideals.
For Lie algebras our categorical use of simple includes a one-dimensional
abelian algebra. Unital rings with identity-preserving maps do not have a zero
object and are not a direct instance of this pointed statement.

\subsection{Reading one cell as the familiar group butterfly}
\label{e:butterfly-dictionary}

Take $A_1\nml A_2\leq G$ and $B_1\nml B_2\leq G$.
Let $r,s$ be the partial equivalence relations on $G$ supplied by
$A_2\twoheadrightarrow A_2/A_1$ and $B_2\twoheadrightarrow B_2/B_1$.
Put
\[
 C=A_2\cap B_2,\qquad U=A_1\cap B_2,\qquad W=A_2\cap B_1.
\]
Applying $s$ and then $r$ to a subset means taking the appropriate
congruence classes inside $B_2$ and then inside $A_2$. In particular,
$sG=B_2$ and $s0=B_1$, so the endpoint identities give
\[
 (rsr)G=rsG=A_1C,\qquad (rsr)0=rs0=A_1W.
\]
Here $0$ denotes the identity subgroup. Interchanging the two relations gives
support $B_1C$ and identity class $B_1U$ for $srs$. Since a group
congruence is determined by its identity class, the two cell quotients are
\[
 Q(rsr)=\frac{A_1C}{A_1W},\qquad
 Q(srs)=\frac{B_1C}{B_1U}.
\]

We can see their common quotient and every normality assertion directly.
The groups $U,W$ are normal in $C$: conjugation by an element of $C$
preserves $A_1$ or $B_1$, respectively, and remains in the other ambient
subgroup. Consequently $UW=WU$ is a normal subgroup of $C$.
The product $A_1C$ is a subgroup because $A_1\nml A_2$ and $C\leq A_2$.
In $A_1C/A_1$, the image of $C$ is the whole group and the image of
$W\nml C$ is normal. Its inverse image is $A_1W$, proving
$A_1W\nml A_1C$. The same argument proves $B_1U\nml B_1C$.

Now use the two homomorphisms
\[
 C\longrightarrow A_1C/A_1W,\quad c\longmapsto c(A_1W),
 \qquad
 C\longrightarrow B_1C/B_1U,\quad c\longmapsto c(B_1U).
\]
Both are onto, since the extra factor $A_1$ or $B_1$ is killed.
The first kernel is $C\cap A_1W=UW$: if $c=aw$ lies in $C$, then
$a=cw^{-1}\in A_1\cap C=U$; the converse inclusion is immediate.
The second kernel is $C\cap B_1U=WU=UW$, by the same calculation.
The effective quotient property therefore gives
\[
\begin{tikzcd}[column sep=large,row sep=large]
 &C\arrow[dl,two heads]\arrow[d,two heads]\arrow[dr,two heads]&\\
 A_1C/A_1W & C/UW\arrow[l,"\sim"']\arrow[r,"\sim"] & B_1C/B_1U .
\end{tikzcd}
\]
These isomorphisms give the usual butterfly, with its maps specified.
They are also the previously constructed cell comparison: for every
element of the correspondence $rs$, its two ends can be represented as
$bc$ and $ac$, with $b\in B_1$, $a\in A_1$ and $c\in C$.
Both quotient maps then use the same class of $c$. The characterization
in Proposition~\ref{r:cell-isomorphism} identifies the maps.
Thus the subgroup products are the concrete supports and identity
classes of one categorical construction.

\subsection{A comparison that changes the order}\label{e:cyclic-six}

In $G=C_6=\mathbb Z/6\mathbb Z$ put $A_1=\{0,3\}$ and
$B_1=\{0,2,4\}$. Reduction modulo $3$ has kernel $A_1$ and reduction
modulo $2$ has kernel $B_1$. Hence
\[
0<A_1<G,\qquad 0<B_1<G
\]
have factors $(C_2,C_3)$ and $(C_3,C_2)$. Every nonzero element of $C_2$
or $C_3$ generates the group, so the factors are simple.
The relation comparison becomes the following concrete rectangle:
\[
\begin{array}{c|cc}
 & B_1/0&G/B_1\\\hline
A_1/0&0&C_2\\
G/A_1&C_3&0
\end{array}
\]
Indeed $A_1\cap B_1=0$ and $A_1+B_1=G$, since $3+4=1$ modulo $6$.
The two nonzero comparisons are induced by $a\mapsto a+B_1$ and
$b\mapsto b+A_1$. Each is onto by the sum equality and injective by the
intersection equality. These are actual isomorphisms, not just equalities
of cardinalities. A refined chain need not contain both $A_1$ and $B_1$;
the matching occurs between its factors.

Normality is always in the next term. In $A_4$, for example,
$1<V_4<A_4$ has both terms normal in the ambient group, but $V_4$ is not
simple. Insert one of its order-two subgroups to obtain a composition series.
That subgroup is normal in the abelian group $V_4$ and is not normal in
$A_4$: conjugation by a three-cycle permutes the three double transpositions.
The lattice of subgroups normal in the ambient group describes a different
filtration problem.

\subsection{Finite basis size, with no dimension in the prerequisites}
\label{e:dimension}

Let $k$ be a field. A list $v_1,\ldots,v_n$ is a basis if it spans the vector
space $V$ and every relation $\sum_i a_i v_i=0$ has all $a_i=0$.
The module construction above provides subspaces, quotient spaces and their
linear maps before any notion of dimension is introduced.

\begin{theorem}[Invariance of finite basis size]\label{e:dimension-theorem}
If $V$ has finite bases $v_1,\ldots,v_n$ and $w_1,\ldots,w_m$, then $n=m$.
\end{theorem}
\begin{proof}
Put $V_i=\Span(v_1,\ldots,v_i)$, with $V_0=0$. Each inclusion is strict:
otherwise $v_i$ would be a linear combination of the preceding vectors,
contradicting independence. The map
\[
k\longrightarrow V_i/V_{i-1},\qquad a\longmapsto av_i+V_{i-1}
\]
is onto by the definition of span. If its image is zero, a relation with
coefficient $a$ on $v_i$ shows $a=0$. It is therefore an isomorphism of
$k$-modules. The module $k$ is simple: a submodule containing $a\ne0$
contains $a^{-1}a=1$ and every scalar multiple of $1$.
The basis thus supplies a composition series with $n$ factors isomorphic
to $k$. The other basis supplies one with $m$ factors. Apply the general
Jordan--H\"older theorem in $k$-modules. If $V=0$, independence forces
both lists empty, giving the same conclusion.
\end{proof}

We may now define $\dim_k V$ to be this invariant basis size. We have not
used bases of arbitrary subspaces, rank--nullity, determinants, exterior
powers, or a pre-existing dimension function. Those are possible later
constructions; they are not concealed inputs to this proof.

\subsection{What a factor list leaves unresolved}\label{e:extensions}

Both $C_4$ and $C_2\times C_2$ have factors $C_2,C_2$. Their exact sequences
can be written
\[
0\longrightarrow C_2\xrightarrow{\,a\mapsto2a\,}C_4
\longrightarrow C_2\longrightarrow0,
\qquad
0\longrightarrow C_2\longrightarrow C_2\times C_2
\longrightarrow C_2\longrightarrow0,
\]
where the last two maps in the second sequence are first-factor inclusion
and second-factor projection. The first quotient is reduction modulo $2$.
The displayed kernels and surjections verify exactness directly.
Yet $C_4$ has an element of order four, whereas every element of
$C_2\times C_2$ is killed by multiplication by two. They are not isomorphic.
The general theorem compares the simple pieces. It does not split the
extension joining them. The following parts develop measurements and
constructions that retain more of this attachment information.
